\section{Further comparison of SAE Types}\label{app:sae_types}

\begin{table*}[htbp]
\centering
\caption{Comparison of some sparse autoencoder (SAE) variants.}
\label{tab:sae_types}
\begin{tabular}{p{2.5cm}p{12.5cm}}
\toprule
\textbf{Type} & \textbf{Description} \\
\midrule
Standard ReLU \cite{bricken2023monosemanticity}
  & Applies an $L_1$ penalty to hidden activations using a standard ReLU activation function to encourage sparsity. \\
\addlinespace
Gated \cite{rajamanoharan2024improving}
  & Decouples feature detection from magnitude estimation using a gated mechanism to prevent the shrinkage issues caused by standard $L_1$ penalties. \\
\addlinespace
JumpReLU \cite{rajamanoharan2024jumping}
  & Uses a discontinuous step function at a learned per-feature threshold, approximating an $L_0$ penalty through straight-through estimation. Notably used in \citet{lieberum2024gemma}. \\
\addlinespace
BatchTopK \cite{bussmann2024batchtopk}
  & Retains top $n \times k$ activations over a batch of $n$ inputs and sets the rest to zero, where $k$ is a tuned hyperparameter. This forces an average of $k$ latent activations per input, so the actual $L_0$ sparsity is approximately $k$. Notably used in \citet{Brixi2026, karvonen2025saebench}. \\
\addlinespace
Matryoshka \cite{bussmann2025learning}
  & Trains nested sub-dictionaries of increasing size that share a single encoder/decoder, producing multiple resolutions of feature granularity in one run. Notably used in \citet{mcdougall2025gemmascope2}. \\
\bottomrule
\end{tabular}
\end{table*}